\documentclass[12pt]{amsart}
%prepared in AMSLaTeX, under LaTeX2e
\addtolength{\oddsidemargin}{-.65in} 
\addtolength{\evensidemargin}{-.65in}
\addtolength{\topmargin}{-.55in}
\addtolength{\textwidth}{1.3in}
\addtolength{\textheight}{.9in}

\renewcommand{\baselinestretch}{1.05}

\usepackage{verbatim} % for "comment" environment

\usepackage{palatino}

\newtheorem*{thm}{Theorem}
\newtheorem*{defn}{Definition}
\newtheorem*{example}{Example}
\newtheorem*{problem}{Problem}
\newtheorem*{remark}{Remark}

\usepackage{fancyvrb,xspace,dsfont}

\usepackage[final]{graphicx}

% macros
\usepackage{amssymb}
\usepackage[dvipsnames]{xcolor}

\usepackage{hyperref}
\hypersetup{pdfauthor={Ed Bueler},
            pdfcreator={pdflatex},
            colorlinks=true,
            citecolor=blue,
            linkcolor=red,
            urlcolor=blue,
            }

\newcommand{\bn}{\mathbf{n}}
\newcommand{\br}{\mathbf{r}}
\newcommand{\bv}{\mathbf{v}}
\newcommand{\bx}{\mathbf{x}}
\newcommand{\by}{\mathbf{y}}

\newcommand{\bJ}{\mathbf{J}}
\newcommand{\bV}{\mathbf{V}}

\newcommand{\cB}{\mathcal{B}}
\newcommand{\cD}{\mathcal{D}}
\newcommand{\cE}{\mathcal{E}}
\newcommand{\cF}{\mathcal{F}}
\newcommand{\cH}{\mathcal{H}}
\newcommand{\cL}{\mathcal{L}}
\newcommand{\cM}{\mathcal{M}}
\newcommand{\cR}{\mathcal{R}}
\newcommand{\cV}{\mathcal{V}}
\newcommand{\cW}{\mathcal{W}}

\newcommand{\CC}{\mathbb{C}}
\newcommand{\NN}{\mathbb{N}}
\newcommand{\RR}{\mathbb{R}}
\newcommand{\ZZ}{\mathbb{Z}}

\renewcommand{\Im}{\mathrm{Im}}
\renewcommand{\Re}{\mathrm{Re}}

\newcommand{\eps}{\epsilon}
\newcommand{\lam}{\lambda}
\newcommand{\lap}{\triangle}

\newcommand{\Div}{\nabla\cdot}
\newcommand{\grad}{\nabla}

\newcommand{\one}{\mathds{1}}

\newcommand{\ip}[2]{\ensuremath{\left<#1,#2\right>}}

\newcommand{\image}{\operatorname{im}}
\newcommand{\onull}{\operatorname{null}}
\newcommand{\rank}{\operatorname{rank}}
\newcommand{\range}{\operatorname{range}}
\newcommand{\trace}{\operatorname{tr}}
\newcommand{\Span}{\operatorname{span}}

\newcommand{\prob}[1]{\bigskip\noindent\textbf{#1.}\quad }

\newcommand{\pts}[1]{(\emph{#1 pts}) }
\newcommand{\epart}[1]{\medskip\noindent\textbf{(#1)}\quad }
\newcommand{\ppart}[1]{\,\textbf{(#1)}\quad }

\newcommand{\ds}{\displaystyle}

\newcommand{\nex}{\medskip\noindent}


\begin{document}
\scriptsize \noindent Math 617 Functional Analysis (Bueler) \hfill \emph{assigned 13 April 2026}
\normalsize\medskip

\Large\centerline{\textbf{Assignment 8}}
\large
\medskip

\centerline{\textbf{Due Friday 24 April at the beginning of class}}
\medskip
\normalsize

\thispagestyle{empty}

\bigskip
\noindent This Assignment is based on the slides, especially those for \href{https://bueler.github.io/fa/assets/slides/S26/week12.pdf}{week 12} and \href{https://bueler.github.io/fa/assets/slides/S26/week12.pdf}{week 13}.

\newcommand{\augment}[1]{\quad\,\strut \begin{minipage}[t]{0.7\textwidth} \emph{#1}\end{minipage}}

\newcommand{\augpart}[1]{\emph{\textbf{(#1)}}}

\medskip
\noindent \textsc{Do the following problems.}


\prob{P23} \emph{Refer to Chapter 4 of the textbook\footnote{K.~Saxe,~\emph{Beginning Functional Analysis}, Springer 2010.} as needed.}

\epart{a} Consider $f\in L^2(-\pi,\pi)$ expanded using the complete ON sequence $\{\frac{1}{\sqrt{2\pi}} e^{inx}\}_{n\in \ZZ}$.  The coefficients $c_n$ in the expansion
	$$f(x) = \frac{1}{\sqrt{2\pi}} \sum_{n=-\infty}^\infty c_n e^{inx}$$
are generally complex, but show that if $f$ is real-valued then
	$$c_{-n} = \overline{c_n}.$$

\epart{b} Again assume that $f$ is real-valued.  Use the substitution $c_n = \frac{\sqrt{2\pi}}{2} (a_n-i b_n)$, with $a_n,b_n\in\RR$, to show the fully-real expansion
    $$f(x) = \frac{a_0}{2} + \sum_{n=1}^\infty a_n \cos(n\theta) + b_n \sin(n\theta).$$
Starting from the formula for $c_n$ as an integral of $f$, derive the simplified, real formulas which compute $a_n$ and $b_n$.

\prob{P24} \emph{Recall the Riesz-Fr\'echet representation theorem:  If $(X,\ip{\cdot}{\cdot})$ is a Hilbert space and if $\ell\in X'$ then there exists a representer $x_\ell\in X$ so that $\ell(x) = \ip{x}{x_\ell}$ for all $x\in X$.}

\medskip \noindent Show that the representer in the Riesz-Fr\'echet representation theorem is unique.


\prob{P25} \emph{A proof of well-posedness of the homogeneous Dirichlet problem for the Poisson equation, using the Riesz-Fr\'echet representation theorem, was shown in the \href{https://bueler.github.io/fa/assets/slides/S26/week13.pdf}{week 13} slides.  Please follow this strategy in part \emph{\textbf{(b)}}.  However, you will \emph{not} need the Poincaré-Friedrichs inequality!}

\epart{a} Suppose $\Omega\subset \RR^d$ is open and bounded, with Lipschitz boundary so that the outward normal vector field $\bn$ is well-defined and the divergence theorem can be applied.  Suppose $\gamma>0$ and $f\in L^2(\Omega)$.  Consider the following strong form of the homogeneous Neumann problem for the Helmholtz equation:
	$$-\grad^2 u + \gamma u = f \quad \text{on } \Omega, \qquad \grad u\cdot \bn = 0 \quad \text{on } \partial \Omega.$$
Derive the weak form, where test functions and the solution are from $H^1(\Omega)$.

\epart{b} Show that the weak form derived in part \textbf{(a)} is well-posed.



\prob{P26} FIXME solve Poisson by series on a square

\prob{P27} FIXME confirm that $\lim_{r\to 1^-} u(r,\theta) = g(\theta)$ with convergence in $L^2$

\prob{P28} FIXME confirm that
\begin{align*}
\grad^2 \phi_{mn} &= - \lambda_{mn}^2 \phi_{mn} \\
\grad^2 \psi_{mn} &= - \lambda_{mn}^2 \psi_{mn}
\end{align*}
use Laplacian in polar, and following equation satisfied by $J_m$

\end{document}
